% Ethics considerations and generative AI usage (IEEE S&P / SaTML requirements).
% In the body of the journal build, in the appendix of the conference build.
\section{Ethics Considerations}
\label{sec:ethics}
\paragraph{Data and human subjects} \bench is derived solely from public, permissively licensed recordings made
by their publishers in lab environments (OTRF Security-Datasets, MIT; Splunk attack\_data, Apache-2.0). They
contain lab host and account names and no personal data; no human subjects were involved and no new data was
collected.

\paragraph{Live systems} No live system was measured, tested or accessed. All investigations ran against a
local runtime serving the recorded graphs; the only external service used was the model provider's API, under
its terms of use.

\paragraph{Tools and attacks} The agent is defensive: it reads telemetry and recommends responses, and takes no
action on systems. The mimicry variant shows that renaming attacker-chosen artefacts can lead an LLM
investigation agent to clear real attacks. Renaming tools and accounts is long-standing intruder tradecraft,
so publishing the result gives attackers little they lack, while defenders learn not to trust names and how to
measure that. Mitigations follow from the analysis: require behaviour-based justification for clearing an
alert, check grounding against identifier types, and evaluate on cue-controlled data. The finding concerns how
investigation agents weigh names, not a flaw in a specific deployed product, so no vulnerability disclosure
process applied.

\paragraph{Availability} Code (MIT), the benchmark and its variants (with the upstream licences), the traces of
every tool call, and the scripts that regenerate every table are released with the paper.

\section{Generative AI Usage}
\label{sec:genai}
\paragraph{As object of study} The agent under study is driven by Claude Opus 5.5 (\texttt{claude-opus-5-5})
through the Anthropic API. The model is closed, and later versions may behave differently; we pin the model
identifier, run three repeats per condition to measure variance, and release the trace of every call so that
results can be inspected without re-running the model.

\paragraph{Footprint} All experiments together (26 runs, 536 investigations) made 3{,}150 model calls, with
0.86\,M output tokens and 19.6\,M prompt tokens read from the provider's cache (2.4\,M written; uncached input
was negligible). Queries were kept low by prompt caching, a per-task call budget, a small benchmark, and offline
rescoring: when the ground truth was revised, earlier runs were rescored from their traces instead of being
repeated. A large model was used because the task requires multi-step tool use over long tool outputs; smaller
models are the planned next comparison. Experiments ran on the provider's hosted infrastructure; local
computation (graph building, scoring, the agent loop) ran on a 4-vCPU cloud server.

\paragraph{In preparing this work} % TODO(author): confirm this statement reflects your actual use before submission.
Generative AI (an LLM-based coding assistant) was used to help implement the software, run and check analyses,
and draft and edit the manuscript. All outputs were inspected and validated by the author, who is responsible
for the content, including the accuracy of every number (each is generated from the released traces) and the
completeness of the literature review.
